\section{Related Work}\label{sec:related}

本节围绕与本文相关的三个方向展开综述：LLM 智能体框架（\S\ref{sec:llm-agents}）、意图识别与对话系统（\S\ref{sec:intent}）以及视觉感知与桌面智能体（\S\ref{sec:visual}）。

\subsection{LLM-based Agent Frameworks}\label{sec:llm-agents}

近年来，基于大语言模型（LLM）的智能体框架迅速发展，推动了从单轮问答到多步任务执行的范式转变。LangChain~\cite{chase2022langchain} 是最早被广泛采用的 LLM 应用框架之一，其核心设计理念是通过"链（Chain）"将提示模板、LLM 调用和外部工具串联为线性流水线。这种链式结构简单直观，但在需要条件分支、循环重试或多步状态管理的复杂任务中表达能力有限。为此，LangChain 团队进一步提出了 LangGraph~\cite{langchain2024langgraph}，引入有向图状态机模型：任务流程被建模为由节点（node）和边（edge）组成的图，节点执行计算并在共享状态上读写，边支持条件路由，从而实现了比线性链更灵活的控制流编排。LangGraph 的设计思想与本文的图工作流引擎直接相关——本文在其基础上构建了 5 个专用子图，分别处理代码生成、文件操作和故障修复等场景。

在多智能体协作方向，AutoGPT~\cite{richards2023autogpt} 率先探索了"自主智能体"的概念，通过让 LLM 自主分解目标、规划步骤和调用工具来完成开放式任务，但其纯自主循环容易陷入低效的重试和目标漂移。MetaGPT~\cite{hong2023metagpt} 提出了另一种思路：将软件工程中的标准化操作流程（SOP）编码为多角色协作协议，为不同智能体分配产品经理、架构师、工程师等角色，通过发布-订阅的消息机制实现结构化通信，在软件开发基准测试上取得了优于 AutoGPT 的表现。CAMEL~\cite{li2023camel} 则聚焦于两个 LLM 角色之间的"角色扮演对话"，通过指令跟随和相互反馈实现任务协作。Voyager~\cite{wang2023voyager} 在 Minecraft 游戏环境中探索了开放式学习，通过代码作为动作空间、技能库的自动积累和课程生成机制，展示了 LLM 智能体在持续探索任务中的潜力。

这些框架的共同趋势是从简单的提示链走向结构化的状态管理和多步编排。然而，它们大多面向通用任务场景，缺乏针对桌面交互场景的分层设计和视觉感知集成。本文继承了 LangGraph 的图状态机思想，但将其聚焦于桌面智能体的三层解耦架构，并加入了视觉上下文感知能力。

\subsection{Intent Recognition and Dialogue Systems}\label{sec:intent}

意图识别是对话系统的核心组件，其技术演进大致可分为三个阶段。第一阶段以规则和关键词方法为主：早期的任务型对话系统如 Rasa~\cite{bocklisch2017rasa} 采用基于模板的意图分类和槽位填充，通过预定义的关键词集合和正则表达式匹配用户输入。这类方法延迟极低且可解释性强，但对表述变化的鲁棒性差，需要大量人工维护规则。

第二阶段以 BERT~\cite{devlin2019bert} 等预训练语言模型为代表，推动了基于神经网络的意图分类。BERT 通过在大规模语料上预训练获得的上下文表征能力，显著提升了对多样化表述的理解。后续工作如 JointBERT 将意图识别与槽位填充统一建模为序列标注任务~\cite{chen2019jointbert}，在多个基准数据集上超越了规则方法。然而，神经网络方法需要大量标注数据，且推理延迟高于规则方法，在对响应速度要求极高的场景中存在局限。

第三阶段利用 LLM 的零样本和少样本能力进行意图分类。GPT-3~\cite{brown2020gpt3} 和后续的指令微调模型展示了无需专门训练即可理解新意图类别的能力。近期研究进一步探索了将 LLM 意图分类与传统方法结合的混合策略：例如用规则方法快速处理明确的输入，将模糊或低置信度的输入交给 LLM 处理~\cite{zhang2022cot}。这种分层处理的思想与本文的路由设计有相似之处，但本文采用的关键词组合门控方案更加轻量——不依赖置信度评分，而是通过多语义类别的布尔共现判定来平衡精度与延迟，且完全不引入额外的 LLM 调用开销。

\subsection{Visual Perception and Desktop Agents}\label{sec:visual}

视觉感知技术为桌面智能体提供了超越文本的环境理解能力。在目标检测领域，YOLO 系列经历了从 YOLOv3~\cite{redmon2018yolov3} 的 anchor-based 设计到 YOLOv5 的工程化优化，再到 YOLOv8~\cite{jocher2023yolov8} 的 anchor-free 架构演进。YOLOv8 采用 C2f 模块替代了 YOLOv5 的 C3 模块，引入解耦头将分类与回归任务分离，并使用 Task-Aligned Assigner 进行动态标签分配，在精度和速度之间取得了更好的平衡。本文使用的微调 YOLOv8 模型正是基于这一架构，针对桌面 UI 元素（窗口、标签页、菜单栏等 6 类）进行了领域适配。

在屏幕理解方向，Screen2Words~\cite{bai2021screen2words} 首次将移动端 UI 截图编码为自然语言描述，为屏幕内容的语义理解奠定了基础。UIBert~\cite{bai2022uibert} 通过多模态预训练（视觉、文本、结构信息的联合编码）提升了对 UI 元素及其关系的理解能力。Ferret-UI~\cite{you2024ferretui} 进一步将大语言模型与屏幕理解结合，实现了对 UI 元素的精确定位和自然语言引用。这些工作推动了屏幕从"像素集合"向"可理解的交互界面"的转变。

在桌面智能体系统方面，Claude Computer Use~\cite{anthropic2024computeruse} 将多模态 LLM 的视觉理解能力直接应用于桌面操作，通过截屏-理解-点击的闭环实现通用的桌面任务自动化，但依赖云端大模型和 GPU 推理。UFO~\cite{zhang2024ufo} 提出了面向 Windows 应用的 UI-Focused Agent，通过双智能体架构（AppAgent 和 ActAgent）分别负责应用选择和操作执行。OS-Copilot~\cite{wu2024oscopilot} 则探索了操作系统级别的通用智能体，能够跨应用协调完成复杂任务。这些系统展示了视觉感知在桌面智能体中的巨大潜力，但均依赖大规模云端模型，难以在资源受限的桌面端实时运行。

本文与上述工作的区别在于：使用轻量级的微调 YOLOv8 模型（12 MB，CPU 推理）实现桌面 UI 元素检测，不追求精细的 UI 语义理解，而是将检测结果映射为粗粒度的活动画像（browsing、managing files 等），作为角色对话的上下文输入。这种设计在保持实时性的同时，为桌面智能体提供了有效的场景感知能力。

\medskip
综合以上三个方向，现有工作在 LLM Agent 框架、意图识别和视觉感知方面各自取得了显著进展，但缺乏一个将多策略意图识别、视觉感知和图工作流在工程层面统一集成并能在桌面端实时运行的轻量级智能体系统。LangGraph 等框架提供了强大的图编排能力但缺少视觉感知；YOLOv8 和屏幕理解技术提供了视觉基础但未与 Agent 工作流深度整合；意图识别的混合策略已有探索但未聚焦于桌面场景的低延迟需求。本文填补了这一工程空白，提出一种三层解耦的桌面智能体架构，在路由层采用关键词组合门控实现零延迟意图分类，在视觉层使用微调 YOLOv8 实现实时活动感知，在工作引擎层基于 LangGraph 实现 5 个子图的图工作流编排，并给出了完整的工程实现与原型验证。
